%% kapa: exam=1 round=22 year=2011 period=1 form=A subject=economics subject_ko=경제원론 seq=8 number=48
%% source: Q-Net 감정평가사 기출문제 (cst00309) artlSeq=5036790 file=extracted/1008/5036790/2011년 제22회 감정평가사 1차시험 기출문제/1교시.hwp
%% conversion: kordoc 4.24.51 -> tools/gen1.py (auto); review=proofread-sonnet; flags=none
\begin{question}{48}
구매력평가설에 관한 설명으로 옳지 않은 것은?
\begin{choices}
  \item 구매력평가설에 의하면 일물일가의 법칙이 성립될 수 있도록 환율이 결정된다.
  \item 절대적 구매력평가설에 의하면 국내 인플레이션율과 해외 인플레이션율은 항상 같다.
  \item 절대적 구매력평가설이 성립하면 실질환율이 1이 된다.
  \item 무역장벽이 높을수록 구매력평가설의 현실 설명력은 감소한다.
  \item 비교역재(non-tradable goods)의 존재가 구매력평가설의 현실 설명력을 떨어뜨리는 요인이 된다.
\end{choices}
\end{question}
\answer{2}
